\documentclass{oxmathproblems}

\printanswers

\course{Understanding Analysis, 2nd Edition}
\sheettitle{Chapter 4 - Section 4.5 - The Intermediate Value Theorem} 

\begin{document}
\begin{questions}

%------------------------- Problem 1 -------------------------
\miquestion Show how the IVT follows as a corollary to Theorem 4.5.2.
\begin{solution}
  Theorem 4.5.2: continuous function preserves connected sets.

  Definition of connected sets: A set $E \subseteq \textbf{R}$ is connected iff whenever $a < c < b$ with $a, b \in E$, then $c \in E$ as well.

  In \textbf{R}, connected sets coincide precisely with the collection of intervals (page 104).
  
  Then IVT follows pretty much by definition. A closed interval $[a, b]$ is connected. Let $f$ be continuous, then $f([a, b])$ is connected. By definition of connectedness,
  $f$ contains everything in $[f(a), f(b)]$.
\end{solution}

%------------------------- Problem 2 -------------------------
\miquestion Example or prove impossible.
\begin{enumerate}[label=(\alph*)]
  \item A continuous function defined on an open interval with range equal to a closed interval.
  \item A continuous function defined on a closed interval with range equal to an open interval.
  \item A continuous function defined on an open interval with range equal to an unbounded closed set different from \textbf{R}.
  \item A continuous function defined on all of \textbf{R} with range equal to \textbf{Q}.
\end{enumerate}

\begin{solution}
  \begin{enumerate}[label=(\alph*)]
    \item $\sin(1/x)$ is continuous on $(0, 1)$ with range $[-1, 1]$.
    \item For unbounded interval, yes. $f(x)=x$ is continuous on \textbf{R} and its range is also \textbf{R} (recall that \textbf{R} is both open and closed).
    
    For unbounded interval, impossible. A closed bounded interval is compact and continuous functions preserve compactness (Theorem 4.4.1).
    \item With unbounded interval: $x^2$ is continuous on \textbf{R} with range $[0, \infty)$.
    
    With bounded interval: $\tan^2{x}$ is continuous on $(\frac{\pi}{2}, \frac{\pi}{2})$ with range $[0, \infty)$.

    \item Impossible. We'll prove it by contradiction. Suppose there exists such a function. Since the range is \textbf{Q}, $f$ is not constant. 
    Choose $x \neq y$ such that $f(x) \neq f(y)$ and wlog $f(x) < f(y)$.
    By IVT, the range of $f$ must contain everything in $[f(x), f(y)]$. But this implies $f$ contains an
    irrational, since between any two rationals there is an irrational, contradicting the assumption that the range is \textbf{Q}.
  \end{enumerate}
\end{solution}

%------------------------- Problem 3 -------------------------
\miquestion A function $f$ is \textit{increasing} on $A$ if $f(x) \leq f(y)$ for all $x < y$ in $A$. Show that if $f$ is increasing on $[a, b]$ and satisfies the
intermediate value property (Definition 4.5.3), then $f$ is continuous on $[a, b]$.
\begin{solution}
  To prove that $f$ is continuous at $a \leq c \leq b$, we need to show that given $\epsilon > 0$, there exists $\delta > 0$ such that $|x-c| < \delta \implies |f(x)-f(c)| < \epsilon$.

  We'll prove left continuity: there exists $\delta_{a}$ such that for $x < c$ and $c-x < \delta_{a}$, $f(c)-f(x) < \epsilon$. The proof for $x > c$ is similar and will result in a $\delta_{b}$. Then
  choose $\delta=\text{min}(\delta_{a}, \delta_{b})$.

  By monotonicity, $f(a) \leq f(c)$ and if $a < x < c$ then $f(a) \leq f(x) \leq f(c)$. By intermediate value property, $f$ takes on all values between $f(a)$ and $f(c)$ so there is an $a < x < c$
  such that $f(c)-f(x) < \epsilon$. Choose $\delta_{a} = c-x$.
  This works for all $x < t < c$ as well because monotonicity ensures that $f(c)-f(t) \leq f(c)-f(x) < \epsilon$.
\end{solution}

%------------------------- Problem 4 -------------------------
\miquestion Let $g$ be continuous on an interval $A$ and let $F$ be the set of points where $g$ fails to be one-to-one; that is,
\[F = \{x \in A : f(x)=f(y) \text{ for some } y \neq x \text{ and } y \in A\}.\]
Show $F$ is either empty or uncountable.

\begin{solution}
  At first I thought the problem statement must be wrong. Take something like $\sin(x)$. Within an interval there's only a finite number of points where $\sin(x)=1$, so it's neither empty nor uncountable. But then I realized
  it's not only 1 that's repeating, the intervals around $\sin(x)=1$ must be repeating as well and an interval is uncountable. That's what we'll prove; $F$ is either empty or contains repeating intervals of output values.

  If $g$ never fails to be one-to-one then $F$ is trivally empty. Otherwise there exists $x \neq y$ such that $g(x)=g(y)$.
  
  If for all $x < c < y$, $g(c)=g(x)=g(y)$, then $[x, y] \subseteq F$ so $F$ is uncountable.

  Otherwise, let $g(c) \neq g(x)$ and wlog $g(c) > g(x)$. By IVT there exists a subset $S_{1} \subseteq (x, c)$ such that $g(S_{1}) = (g(x), g(c))$. Since $g(S_{1})$ is an interval, $S_{1}$ must be uncountable.
  
  Similarly, there exists a subset $S_{2} \subseteq (c, y)$ such that $g(S_{2})=(g(y), g(c))$.
  
  Since $g(x)=g(y)$ then $g(S_{1})=g(S_{2})$, hence $S_{1}, S_{2} \subseteq F$ and $F$ is uncountable.  Note that $S_{1} \cap S_{2} = \emptyset$ so they're different set of points
  and can be included in $F$.
\end{solution}

%------------------------- Problem 5 -------------------------
\miquestion
\begin{enumerate}[label=(\alph*)]
  \item Finish the proof of the Intermediate Value Theorem using the Axiom of Completeness started previously.
  \item Finish the proof of the Intermediate Value Theorem using the Nested Interval Property started previously.
\end{enumerate}

\begin{solution}
  \begin{enumerate}[label=(\alph*)]
    \item The partial proof from the book is as follow. Let $f$ be continuous and consider the special case $f(a) < 0 < f(b)$. We'll show that $f(c)=0$ for some $c \in (a, b)$. Let
    \[K = \{x \in [a, b] : f(x) \leq 0\}.\]
    Notice that $K$ is bounded above by $b$ and $a \in K$ so $K$ is not empty. Thus we may appeal to AoC to assert that $s=\text{sup }K$ exists.

    We'll first prove that $f(s) \leq 0$. Since $s$ is a supremum, $s$ is a limit point of some sequence $(x_{n}) \in K$. By continuity, $f(x_{n}) \to f(s)$. Since
    $f(x_{n}) \leq 0$, then $f(s) \leq 0$.

    Next we need to prove that $f(s) \geq 0$. The idea is that since $s$ is a supremum, it's sort of the "rightmost" point where $f$ is $\leq 0$. But if $f(s) < 0$,
    since everything to the right of $s$ is $> 0$, there'd be a discontinuity; a jump that skips 0. Or in other words, $f(s)$ is not continuous from the right side.

    Formally, we prove by contradiction. First we note that since $s = \text{sup } K$, if $x > s$ then $x \notin K$ hence $f(x) > 0$.
    
    Now let $\epsilon = |f(s)|$. By continuity, there exists a $\delta < (b-s)$ that works (the bound $(b-s)$ is just a trivial detail to ensure
    we stay within $[a, b]$). Choose $x$ such that $x-s < \delta$. Since $f$ is defined on $[a, b]$, $f(x)$ is defined.
    
    We've asserted that $f(x) > 0$. But since $f(s) < 0$, we have $f(x)-f(s) \geq |0-f(s)| = \epsilon$, contradicting continuity of $f$.
  \end{enumerate}
\end{solution}

\end{questions}
\end{document}
